%! Matrices and Column Spaces — Unit 1, Lesson 1.3 — Homework (Answer Key)
\documentclass[10pt]{article}
\usepackage{linalg-article}
\usepackage{linalg-key}   % pulls in -boxes; do NOT also load -boxes
\newcommand{\TallMath}[1]{$\displaystyle #1\rule[-1.4em]{0pt}{3.2em}$}
\newcommand{\xx}{\mathbf{x}}
\newcommand{\bb}{\mathbf{b}}

\begin{document}

\pageheader{Unit 1, Lesson 1.3}{Homework --- Answer Key}

\vspace{0.12in}

\begin{notesbox}{Practice}
For Problems 1--2, let $A = \begin{bmatrix} 1 & 2 \\ 3 & 1 \end{bmatrix}$. Show your steps.

\begin{enumerate}[leftmargin=1.9em, itemsep=10pt]
  \item Write the two columns of $A$, then compute each product as a combination of the columns.
    \begin{enumerate}[label=(\alph*), leftmargin=1.8em, itemsep=4pt]
      \item $A\begin{bsmallmatrix} 1 \\ 0 \end{bsmallmatrix}$ \hfill
      \item $A\begin{bsmallmatrix} 0 \\ 1 \end{bsmallmatrix}$ \hfill
      \item $A\begin{bsmallmatrix} 2 \\ 1 \end{bsmallmatrix}$ \hfill
      \item $A\begin{bsmallmatrix} 1 \\ 3 \end{bsmallmatrix}$
    \end{enumerate}
\begin{work}
  \mathbf{a}_1 &= \begin{bsmallmatrix} 1\\3 \end{bsmallmatrix} \\
  \mathbf{a}_2 &= \begin{bsmallmatrix} 2\\1 \end{bsmallmatrix} \\
  \text{(a)}\ 1\mathbf{a}_1 + 0\mathbf{a}_2 &= \begin{bsmallmatrix} 1\\3 \end{bsmallmatrix} \\
  \text{(b)}\ 0\mathbf{a}_1 + 1\mathbf{a}_2 &= \begin{bsmallmatrix} 2\\1 \end{bsmallmatrix} \\
  \text{(c)}\ 2\mathbf{a}_1 + 1\mathbf{a}_2 &= \begin{bsmallmatrix} 4\\7 \end{bsmallmatrix} \\
  \text{(d)}\ 1\mathbf{a}_1 + 3\mathbf{a}_2 &= \begin{bsmallmatrix} 7\\6 \end{bsmallmatrix}
\end{work}
  \item Is the column space $C(A)$ a line or the whole plane? Explain how you can tell from the
        columns. \ansline{Whole plane --- $\begin{bsmallmatrix} 1\\3 \end{bsmallmatrix}$ and $\begin{bsmallmatrix} 2\\1 \end{bsmallmatrix}$ are not multiples of each other (independent), so their combinations reach every point.}
  \item For each matrix, decide whether its column space is a \textbf{line} or the whole
        \textbf{plane}. Give a one-line reason.
    \begin{enumerate}[label=(\alph*), leftmargin=1.8em, itemsep=4pt]
      \item $\begin{bsmallmatrix} 1 & 2 \\ 2 & 4 \end{bsmallmatrix}$ \hfill
      \item $\begin{bsmallmatrix} 1 & 0 \\ 0 & 1 \end{bsmallmatrix}$ \hfill
      \item $\begin{bsmallmatrix} 3 & 1 \\ 6 & 2 \end{bsmallmatrix}$
    \end{enumerate}
    \ansline{(a) Line --- $\begin{bsmallmatrix} 2\\4 \end{bsmallmatrix}=2\begin{bsmallmatrix} 1\\2 \end{bsmallmatrix}$ (parallel).
    \ (b) Plane --- $\begin{bsmallmatrix} 1\\0 \end{bsmallmatrix}$ and $\begin{bsmallmatrix} 0\\1 \end{bsmallmatrix}$ are independent.
    \ (c) Line --- $\begin{bsmallmatrix} 3\\6 \end{bsmallmatrix}=3\begin{bsmallmatrix} 1\\2 \end{bsmallmatrix}$ (parallel).}
  \item Use the columns to decide whether each target $\bb$ is in the column space.
    \begin{enumerate}[label=(\alph*), leftmargin=1.8em, itemsep=4pt]
      \item Is $\begin{bsmallmatrix} 4 \\ 8 \end{bsmallmatrix}$ in $C\!\left(\begin{bsmallmatrix} 1 & 3 \\ 2 & 6 \end{bsmallmatrix}\right)$? \hfill
      \item Is $\begin{bsmallmatrix} 1 \\ 0 \end{bsmallmatrix}$ in $C\!\left(\begin{bsmallmatrix} 1 & 2 \\ 2 & 4 \end{bsmallmatrix}\right)$?
    \end{enumerate}
    \ansline{Both column spaces are the line of multiples of $\begin{bsmallmatrix} 1\\2 \end{bsmallmatrix}$.
    \ (a) $\begin{bsmallmatrix} 4\\8 \end{bsmallmatrix}=4\begin{bsmallmatrix} 1\\2 \end{bsmallmatrix}$ --- yes, in $C$.
    \ (b) $\begin{bsmallmatrix} 1\\0 \end{bsmallmatrix}$ is not a multiple of $\begin{bsmallmatrix} 1\\2 \end{bsmallmatrix}$ --- no.}
  \tcbbreak
  \item A workshop builds two kits as columns of $A=\begin{bmatrix} 3 & 1 \\ 1 & 2 \end{bmatrix}$,
        each listing $\begin{bsmallmatrix} \text{bolts} \\ \text{brackets} \end{bsmallmatrix}$ per kit.
    \begin{enumerate}[label=(\alph*), leftmargin=1.8em, itemsep=5pt]
      \item Compute $A\begin{bsmallmatrix} 2 \\ 1 \end{bsmallmatrix}$ and say what it means in
            bolts and brackets.
\begin{work}
  A\begin{bsmallmatrix} 2\\1 \end{bsmallmatrix}
      &= 2\begin{bsmallmatrix} 3\\1 \end{bsmallmatrix} + 1\begin{bsmallmatrix} 1\\2 \end{bsmallmatrix} \\
      &= \begin{bsmallmatrix} 7\\4 \end{bsmallmatrix}
\end{work}
            \ansline{2 of kit~1 and 1 of kit~2 use 7 bolts and 4 brackets.}
      \item Can the workshop hit \emph{any} target of bolts and brackets by choosing how many of
            each kit? Explain using the columns. \ansline{Yes --- $\begin{bsmallmatrix} 3\\1 \end{bsmallmatrix}$ and $\begin{bsmallmatrix} 1\\2 \end{bsmallmatrix}$ are independent, so $C(A)$ is the whole plane and every target is reachable.}
    \end{enumerate}
\end{enumerate}
\end{notesbox}

\begin{extensionbox}
A factory's two product bundles are the columns of $A=\begin{bmatrix} 2 & 4 \\ 3 & 6 \end{bmatrix}$,
each listing $\begin{bsmallmatrix} \text{labor hours} \\ \text{material units} \end{bsmallmatrix}$.
The manager wants a production plan using $x_1$ of bundle~1 and $x_2$ of bundle~2 that totals
exactly $\begin{bsmallmatrix} 5 \\ 6 \end{bsmallmatrix}$. Explain, using the column space, why
\emph{no} plan can hit this target --- and describe every total the factory \emph{can} hit.
\ansline{Column~2 $=2\times$ column~1, so $C(A)$ is only the line of multiples of $\begin{bsmallmatrix} 2\\3 \end{bsmallmatrix}$.
$\begin{bsmallmatrix} 5\\6 \end{bsmallmatrix}$ is not a multiple of $\begin{bsmallmatrix} 2\\3 \end{bsmallmatrix}$ ($5=2t\Rightarrow t=2.5$ but $6=3t\Rightarrow t=2$), so no plan works. The factory can hit exactly the totals $\begin{bsmallmatrix} 2t\\3t \end{bsmallmatrix}$ --- labor and material in a fixed $2\!:\!3$ ratio.}
\end{extensionbox}

\begin{spiralbox}
\textbf{Coming next --- Lesson 1.4: Matrix Multiplication and $A=CR$.} You can now read a matrix as
columns and see its column space. Next we build \emph{every} column of a matrix from a few
\emph{independent} columns --- the factorization $A=CR$.
\end{spiralbox}

\end{document}
